% --- FONT E CODIFICA ---
\usepackage[utf8]{inputenc}
\usepackage[T1]{fontenc}
\usepackage[italian]{babel}

% Libertinus: Leggibile, moderno e perfetto per la matematica
\usepackage{libertinus} 
% Fira Mono: Il font per il codice (scalato leggermente per matchare il testo)
\usepackage[scaled=0.85]{beramono}
% --- MATEMATICA AVANZATA ---
\usepackage{amsmath, amssymb, amsthm}
\usepackage{mathtools} % Estensione di amsmath
\usepackage{float}
\usepackage{ulem}
\usepackage{tcolorbox}
\usepackage{tabularx}

% --- COLORI ---
\usepackage{xcolor}
\definecolor{NoteRed}{RGB}{180, 0, 0}
\definecolor{CodeBack}{RGB}{245, 245, 245} % Grigio chiarissimo per lo sfondo del codice
\definecolor{CodeGreen}{RGB}{0, 120, 0}
\definecolor{CodeGray}{RGB}{128, 128, 128}

% --- CONFIGURAZIONE CODICE (Listings) ---
\usepackage{listings}
\lstset{
    backgroundcolor=\color{CodeBack},
    basicstyle=\ttfamily\small,
    breakatwhitespace=false,
    breaklines=true,
    captionpos=b,
    commentstyle=\color{CodeGreen},
    keywordstyle=\color{NoteRed}\bfseries, % Parole chiave in rosso come il tuo tema
    stringstyle=\color{orange},
    numbers=left,
    numberstyle=\tiny\color{CodeGray},
    numbersep=8pt,
    frame=single,
    rulecolor=\color{lightgray},
    showstringspaces=false,
    language=C++, % Linguaggio di default (puoi cambiarlo nel main)
}

% --- RESTO DELLE TUE CONFIGURAZIONI ---
\usepackage{geometry}
\geometry{a4paper, margin=2cm}
\usepackage{parskip}
\usepackage{titlesec}

% Hack Grassetto Rosso
\let\oldtextbf\textbf
\renewcommand{\textbf}[1]{\textcolor{NoteRed}{\oldtextbf{#1}}}

% Titoli Rossi
\titleformat{\chapter}[display]{\normalfont\Huge\bfseries\color{NoteRed}}{\chaptertitlename\ \thechapter}{10pt}{\Huge}
\titleformat{\section}{\normalfont\Large\bfseries\color{NoteRed}}{\thesection\quad}{0pt}{}
\titleformat{\subsection}{\normalfont\large\bfseries\color{NoteRed}}{\thesubsection\quad}{5pt}{}

\usepackage[colorlinks=true, allcolors=NoteRed]{hyperref}